\chapter*{Presentación}
\addcontentsline{toc}{chapter}{Presentación}
Este manual expone la responsabilidad extracontractual del Estado en Colombia: sus antecedentes en el derecho civil y en la historia del derecho administrativo, la cláusula general del artículo 90 de la Constitución, los elementos de la responsabilidad (daño y nexo causal), los títulos de imputación, la responsabilidad por la actividad judicial y por el hecho del legislador, el medio de control de reparación directa, las reformas de la Ley 2195 de 2022 y la responsabilidad patrimonial del agente estatal mediante la acción de repetición.

El contenido se elaboró exclusivamente a partir del material del curso de Responsabilidad extracontractual del Estado (diapositivas, apuntes y exposiciones de clase). Cada tema se desarrolla con definiciones, clasificaciones, normas, jurisprudencia y ejemplos, e incluye las principales críticas y debates sobre la materia.

\textbf{Convenciones.}
\begin{definicion}Definiciones de figuras y conceptos.\end{definicion}
\begin{ejemplo}Ejemplos y casos, reales o hipotéticos.\end{ejemplo}
\begin{comentario}Comentarios, críticas, debates y remisiones entre capítulos.\end{comentario}
\begin{cita}Citas textuales de normas, sentencias y doctrina.\end{cita}

\textbf{Abreviaturas.} C.C.: Código Civil. C.P.: Constitución Política. CPACA: Código de Procedimiento Administrativo y de lo Contencioso Administrativo (Ley 1437 de 2011). CCA: Código Contencioso Administrativo (anterior). CGP: Código General del Proceso. CE: Consejo de Estado. S3: Sección Tercera. CSJ: Corte Suprema de Justicia. Corte IDH / CIDH: Corte Interamericana de Derechos Humanos. SMLMV / SMMLV: salarios mínimos legales mensuales vigentes. ANDJE: Agencia Nacional de Defensa Jurídica del Estado. ANTV: Autoridad Nacional de Televisión.

